\chapter*{Introduction}
\addcontentsline{toc}{chapter}{Introduction}  

\begin{chapquote}{\textcite{ormesson_cest_2010}, \textit{C'est une chose étrange à la fin que le monde}}
	 ``La science d'aujourd'hui détruit l'ignorance d'hier et elle fera figure d'ignorance au regard de la science de demain. Dans le coeur des hommes il y a un élan vers autre chose dont la clé secrète est ailleurs.''
\end{chapquote}%Jean d'Ormesson

%
C'est bien souvent en observant son environnement, et en trouvant ses muses dans la nature que l'esprit humain trouve l'inspiration. Ainsi Sénèque disait \textit{omnis ars naturae imitatio est}\footnote{Littéralement "Tout art est une immitation de la nature"} \parencite[Lettre~LXV]{seneque_le_jeune_lettres_1914}, repris plus d'un millénaire après par Pablo Picasso, \textit{la peinture ce n'est pas copier la nature mais c'est apprendre à travailler comme elle}\footnote{Citation possiblement apocryphe}. Ainsi, ces deux citations reflètent les deux approches différentes que sont le bio-mimétisme et la bio-inspiration. La première étant souvent associée à la recherche d'une compréhension ou de la pure reproduction d'un phénomène, et la seconde plutôt de la recherche d'une même finalité.

Cette ambition mimétique est assez bien incarnée par certains chercheur japonais, comme Hiroshi Ishiguro du laboratoire de robotique intelligente de l'université d'Osaka, qui a dirigé des équipes pour développer et concevoir des robots dont l'apparence et les mimiques s'approchent autant que possible d'humain \parencite{hiroshi_how_2020}. La démarche est ici presque artistique, même si le chercheur avance une volonté de faire progression la compréhension des interactions sociales humaines et de manière plus générale des siences cognitives. Reste que ces robots sont incapables d'évoluer de manière autonome dans des environnements humains.

Le système neuro-musculo-squelettique (SNMS) humain est en effet d'une telle complexité qu'il limite les approches mimétiques. La compréhension des performances de la motricité humaine est d'ailleurs l'un des axe d'étude sur lequel se positionne l'unité de recherche CIAMS\footnote{Complexité, Innovation, Activités Motrices et Sportives} de l'université Paris Saclay. Car malgré le chantier immense que représente le SNMS, il n'en reste pas moins une grande source d'inspiration, notamment pour la robotique. Les différences entre la biologie humaine et l'électro-mécanique de la robotique pousse plutôt à essayer d'approcher les performances humaines en s'inspirant de leurs stratégies variées.

%
En effet, les robots de plus en plus présents dans les environnements anthropisés doivent interagir avec des outils pensés pour l'Homme, mais aussi avec des êtres humains. Dans ce contexte, une meilleure compréhension des mouvements humains est essentielle pour ces interactions. Par ailleurs, les capacités dextres humaines ainsi que la perception de l'environnement restent aujourd'hui un grand défi pour la robotique. Ces défis se placent aussi bien sur le plan mécanique avec les différences entre les complexes musculo-squelettiques et les actionneurs robotisés que sur le plan de l'automatique avec le contrôle neural face aux lois de commande en robotique. C'est sur ce second axe que se positionnent certains chercheurs du L2S\footnote{Laboratoire des Signaux et Systèmes}.

%
% C'est ainsi qu'une collaboration entre enseignants-chercheurs des deux structures a pu naître.
La modélisation du membre supérieur contribue d’une part à la compréhension du mouvement humain et à son analyse à l’aide d’outils de l’automatique, et d’autre part au développement de stratégies de commande bio-inspirées innovantes pour des systèmes robotiques anthropomorphes en interaction avec l’environnement, pour des applications en co-manipulation et assistance robotique. \textcite{avrin_modelisation_2017} a proposé un modèle comportemental innovant fondé sur un oscillateur neuronal bio-inspiré, capable de reproduire des motifs de comportement humain, dans le cadre d'une tâche de jonglerie. \textcolor{red}{[expliquer pourquoi cette tâche en qlq mots ??]}

%
Dans la perspective d’un transfert des modèles de comportement ainsi mis en évidence vers des applications robotiques réalistes en interaction dynamique avec des environnements physiques, le modèle développé doit être complété par des stratégies du contrôle des efforts employés pour réaliser le mouvement. Par ailleurs, comprendre comment les propriétés dynamiques du bras sont régulées au cours du mouvement pour préparer des phases d’impact ou de contact (\cite{tsuji_bio-mimetic_2008}) est également un enjeu dans le champ des sciences du mouvement humain. Des mesures d'impédance pourraient permettre d'éclaircir cet aspect et la transposition des connaissances acquises pour proposer une structure de commande innovante pour la robotique.%, avec des analyses de robustesse essayant de vérifier l'intérêt des variations d'impédance pour l'Homme et le robot. 

C'est donc en s'inscrivant dans ces objectifs qu'un banc expérimental a été mis en place pour mener à bien des mesures sur plusieurs personnes pendant une interaction robotique avec des retours haptiques, nécessitant un contrôle actif de participant pour réaliser une tâche de jonglerie dans un environnement virtuel parfaitement maîtrisé.

Les travaux de recherche de cette thèse sont retranscrits dans ce manuscrit qui se découpe en cinq chapitres. Le \cref{chap1} présente le contexte de la robotique industrielle et de service en se focalisant sur les défis qui se posent aujourd'hui pour les interactions physiques et les partages de zone de travail avec des humains. Pour mieux maîtriser et qualifier les interactions physiques avec des humains, une revue de la littérature de la modélisation du comportement viscoélastiques des membres humains, et plus précisément du membre supérieure est livrée. Le concept clée d'impédance mécanique, qui permet de manière très simple de reproduire des comportements humains, est alors introduit. Il faut bien le différencier de l'appelation de contrôle en impédance, aussi abordé dans cette thèse, qui consiste notamment à simuler un comportement en impédance mécanique pour un robot.

Le \cref{chap2} délivre une courte revue de la littérature concernant les différents types de contrôles en impédance qui permettent aux robots d'interargir avec leur environnement en contrôlant indirectement la force et la position par la simulation de comportements viscoélastiques. L'utilisation d'actionneurs innovants permettant de conférer des dynamiques viscoélastiques aux articulations des robots est aussi abordée. Le contrôle mis en place pour le banc expérimental conçu est ensuite détaillé, avec les différents outils de l'automatique utilisés et la méthode de réglage des différents contrôleurs impliqués.

Un rapide panorama des différentes techniques et méthodes permettant d'estimer les propriétés viscoélastiques et de manière plus spécifique d'identifier les paramètres en impédance cartésienne est élaboré au \cref{chap3}. Les aspects concernant les méthodes de perturbation ainsi que d'identification sont abordés pour présenter leurs avantages et limites. Les méthodes proposées dans le cadre de cette thèse sont ensuite explicités, avec un effort particulier placé sur les estimations de trajectoires en force et position nécessaire pour les identifications d'impédance dans le cadre de tâches dynamiques. La méthode d'identification proposée estfinalement évaluée dans une simulation impliquant des signaux d'entrées expérimentaux.

Le cadre expérimental est enfin présenté dans le \cref{chap4}. Le matériel, les protocoles, la cohorte recrutée, ainsi que les hypothèses sont donnés dans une première partie afin de bien comprendre la méthodologie expérimentale utilisée. Puis dans une seconde partie, les données obtenues sont exploitées à l'aide de la méthodologie présenté dans le chapitre précédent, et les résultats sont divulgués grâce à des outils statistiques classiques. Les différents sources de potentielle variabilité sont analysés aussi bien pour les performances de la tâche que pour les données biomécaniques, et sont ensuites discutés en fonction de leur significativité.

Afin de compléter les études conduites, une analyse de stabilité et robustesse est conduite avec un modèle linéarisé du couplage entre le robot et un environnement modélisé en impédance dans le \cref{chap5}. La modélisation dynamique du robot est détaillée, afin que les limites de cette analyse soient claires. Le compromis entre stabilité et transparence mécanique du robot est sondé en utilisant les résultats expérimentaux du \cref{chap4}.

% introduire le chapitre 5 % étude de la stabilité



%
%L'ensemble de l'état de l'art sur ces thématiques, aussi bien du côté des sciences du mouvement humain que de la robotique/automatique, mené au début de la thèse, n'a pas pu être inclus dans son intégrité dans ce rapport par souci de brièveté. Les revues suivantes concernant l'impédance humaine et ses lien avec la robotique sont citées à titre indicatif : \cite{mizrahi_mechanical_2015,vanderborght_variable_2013,nikooyan_mass-spring-damper_2011}, et chapitres de livres : \cite{burdet_interaction_2014,tsuji_human_1997}.

%annonce du plan
